\section{Budgets, private coordinates, and compression}
\label{sec:compression}

We first make precise the quantitative conventions and the projection estimates
used below.  All trees in this section are finite.  Bounds are independent of
their number of generations.

\subsection{Quantitative conventions and stopping attempts}

A rational number has height at most $B$ if it can be written as $p/q$ with
$p\in\Z$, $q\in\N$, and $\max(|p|,q)\le B$.  Heights of rational vectors and
matrices are entrywise heights.  For parameters $H_1,\ldots,H_s\ge e$ and
$0<\delta_1,\ldots,\delta_t\le1$, a quasipolynomial bound means a bound of the
form
\[
 \exp\!\left(C_q\left(1+\sum_{a=1}^s\log H_a
                  +\sum_{b=1}^t\log(1/\delta_b)\right)^C\right).
\]
A polylogarithmic bound is the same expression without the outer exponential.
Here and below the degrees denoted by $C$ can be fixed independently of $q$,
of spatial lengths and frequencies, and of the number of scales.  Constants
such as $C_q$ may depend on the fixed choice of $q$ and on the fixed kernel
class.  A bounded number of compositions of quasipolynomial constructions is
again quasipolynomial.  When the number of constructions is itself a parameter,
we keep track of the resulting growth explicitly.  Expressions such as
$D\le\exp(k^C)$ allow an increase of the fixed exponent and a sufficiently
large lower threshold for $k$.

We will use elementary spatial packing in the following form.  If $J$ is a
dyadic interval, $a_I\ge0$, and
\[
 \sum_{I\ni y}a_I\le A \qquad\text{for almost every }y\in J,
\]
where all the intervals in the sum are contained in $J$, then Tonelli's
theorem gives
\[
 \sum_I |I|a_I
   =\int_J\sum_{I\ni y}a_I\,dy\le A|J|.
\]
This is the passage from a bound along a spatial path to a length-weighted
tree bound.

\begin{lemma}[Packing of attempts]
\label{lem:attempt-packing}
Suppose a construction on a root $J$ sends exceptional cells and all their
descendants to new estimates on disjoint subroots.  Suppose that, at every
root $R$, the new roots $R'$ satisfy
\[
 \sum_{R'\text{ child root of }R}|R'|\le\rho |R|,
 \qquad 0\le\rho<1.
\]
Then the sum of the lengths of all roots, including $J$, is at most
$|J|/(1-\rho)$.  In particular, if the contribution assigned to an attempt
from $R$ is at most $A|R|$, with its hypotheses valid at every new root, the
total contribution is at most $A|J|/(1-\rho)$.
\end{lemma}

\begin{proof}
The sum of the root lengths in generation $n$ is at most $\rho^n|J|$, by
induction.  Summing this geometric series proves both assertions.
\end{proof}

We refer to one application of the construction at a root as an
\emph{attempt}.  Several stopping requirements can be imposed together: if
their exceptional root families cost fractions $\rho_1,\ldots,\rho_s$, take
the maximal intervals in their union.  These are disjoint, their total length
is at most $(\sum_a\rho_a)|R|$, and Lemma~\ref{lem:attempt-packing}
applies when this sum
is smaller than one.  A local normalization required by an estimate must hold
at each new root before that estimate is restarted there.

A stopping boundary for the summation does not have to be a boundary for an
auxiliary projection family.  An inherited column list can be continued on
all descendants of a terminal decision cell: retain its columns and refit
their constant coefficients on each finer cell.  This gives genuine nested
projections on the full refining partitions needed for an operator estimate.
If a later attempt uses a new family, the extension of the old family is
still used for the summands assigned to the old attempt.  In particular, a
projection on a partition refining $I$ is an $L^2(I)$ contraction whether or
not the averages of the original input are normalized on the finer cells.

\subsection{Columns and inherited projections}

Fix a finite dyadic tree with root $J$ and a finite label set $\mathcal V$.
The cardinality of $\mathcal V$ will not be bounded.  To each label $v$ assign
a birth cell $J_v\subset J$, a scalar function $\varphi_v$ on $J_v$, and a
number $\epsilon_v>0$.  Work in the value Hilbert space
\[
 \mathcal H=\C\oplus\ell^2(\mathcal V),
\]
and define the column on its birth cell by
\[
 \Phi_v(y)=\bigl(\varphi_v(y),\sqrt{\epsilon_v}\,\mathbf e_v\bigr),
 \qquad y\in J_v.
\]
The first coordinate is called the physical coordinate.  Every label has its
own private direction, even when two scalar functions $\varphi_v$ coincide.
We write $\widetilde\Phi_v=\1_{J_v}\Phi_v$ for the extension by zero to $J$.

Throughout the projection estimates, $D\ge2$ is a common budget such that
\begin{equation*}
 \|\Phi_v\|_{L^\infty(J_v;\mathcal H)}\le D,
 \qquad \sqrt{\epsilon_v}\ge D^{-1},
 \qquad \sum_{v\in\mathcal V}\1_{J_v}(y)\le D
       \quad\text{for almost every }y\in J.
\end{equation*}
The last condition says that at most $D$ labels are born along any spatial
path.  It counts all available labels, including labels not selected in the
projection currently under consideration.  It allows arbitrarily many labels
on disjoint branches.

A mask $\mathcal V^i(I)$ on a cell $I$ is a subset of
$\{v:I\subset J_v\}$.  An inherited mask family satisfies
\[
 \mathcal V^i(I)\subset\mathcal V^i(I')
 \qquad\text{whenever }I'\text{ is a child of }I.
\]
Let $P_I^i$ be the orthogonal projection in $L^2(I;\mathcal H)$ onto
\[
 \left\{\sum_{v\in\mathcal V^i(I)}c_v\Phi_v|_I:c_v\in\C\right\}.
\]
The coefficients in this definition are constant on $I$.  We also regard
$P_I^i$ as an operator on $L^2(J;\mathcal H)$ by restricting its input to $I$
and extending its output by zero.  For a stopping partition $\pi$ of a cell,
$P_\pi^i$ denotes the direct sum of these local projections.  If $\pi'$
refines $\pi$, inheritance implies
\[
 \operatorname{range}(P_\pi^i)
     \subset\operatorname{range}(P_{\pi'}^i).
\]
Indeed, the restriction of an old constant-coefficient combination to a
smaller cell is still an allowed combination there.  Thus these direct-sum
projections are nested, even though their coefficients are refitted on every
cell.  Their successive differences are orthogonal projections with mutually
orthogonal ranges.

The masks and birth lists may have been chosen using an input or previous
stopping decisions.  In each application of an operator estimate we hold the
entire resulting family fixed.

\begin{lemma}[Coefficient and projection bounds]
\label{lem:private-columns}
If $\mathcal A$ is any set of columns available on a cell $I$, then
\[
 D^{-1}\left(\sum_{v\in\mathcal A}|c_v|^2\right)^{1/2}
 \le \left\|\sum_{v\in\mathcal A}c_v\Phi_v(y)\right\|_{\mathcal H}
 \le D^{3/2}\left(\sum_{v\in\mathcal A}|c_v|^2\right)^{1/2}
 \quad\text{for almost every }y\in I.
\]
In particular, the coefficients are unique.  The orthogonal projection
$P_{I,\mathcal A}$ onto their span satisfies
\[
 \|P_{I,\mathcal A}z\|_{L^\infty(I;\mathcal H)}
       \le D^5\fint_I\|z(y)\|_{\mathcal H}\,dy.
\]
For an actual scalar input $F=(f,0)$, write
$P_{I,\mathcal A}F=\sum_{v\in\mathcal A}c_v\Phi_v$.  Then
\[
 \fint_I\bigl(f-(P_{I,\mathcal A}F)_{\phys}\bigr)
                         \overline{\varphi_v}
       =\epsilon_v c_v \qquad(v\in\mathcal A).
\]
\end{lemma}

\begin{proof}
The private component of the sum has squared norm
$\sum_v\epsilon_v|c_v|^2$, giving the lower bound.  The upper bound follows
from Cauchy--Schwarz, since there are at most $D$ columns and each has norm at
most $D$.

Let $E_y:\ell^2(\mathcal A)\to\mathcal H$ be the evaluation map
$E_yc=\sum_vc_v\Phi_v(y)$, and let
\[
 G=\fint_I E_y^*E_y\,dy
\]
be the normalized Gram matrix.  The bounds just proved give
$G\ge D^{-2}\operatorname{Id}$ and $\|E_y\|\le D^{3/2}$.  Therefore the
projection has the integral representation
\[
 P_{I,\mathcal A}z(y)
      =\fint_I E_yG^{-1}E_t^*z(t)\,dt,
 \qquad \|E_yG^{-1}E_t^*\|\le D^5.
\]
This proves the supremum estimate.  Finally, orthogonality to $\Phi_v$ says
that the physical pairing in the last assertion minus $\epsilon_vc_v$ is
zero.
\end{proof}

\subsection{A binary-prefix estimate}

The following elementary Hilbert space fact supplies the dependence on the
number of stopping generations.  In the compression argument that number
will be a power of $D$, so its logarithm is the appropriate cost.
Its proof uses the classical Rademacher--Menshov binary-decomposition
mechanism; see \cite[Section~3, Lemma~1]{MM2011} for a vector-valued
presentation.  The nested projection products and Bessel rows are treated
below.

\begin{lemma}[Products of nested prefixes]
\label{lem:projection-prefixes}
Let $\mathcal X$ be a Hilbert space, and for $1\le i\le m$ let
$P_0^i,\ldots,P_T^i$ be nested orthogonal projections on $\mathcal X$.
Set $Z_n=P_n^m\cdots P_n^1$ and
\[
 L_T=1+\left\lceil\log_2(T+1)\right\rceil.
\]
Suppose $B_\nu:\mathcal X\to\mathcal Y_\nu$ are finitely many bounded linear operators
satisfying $\sum_\nu\|B_\nu h\|^2\le\|h\|^2$.  For any choices
$n(\nu)\in\{0,\ldots,T\}$,
\[
 \left(\sum_\nu\|B_\nu Z_{n(\nu)}g\|^2\right)^{1/2}
       \le L_T^m\|g\|.
\]
Moreover,
\[
 \sum_{n=1}^T\|(Z_n-Z_{n-1})g\|^2
       \le m^2L_T^{2(m-1)}\|g\|^2
       \le m^2L_T^{2m}\|g\|^2.
\]
For $m=0$ the product is the identity and the first estimate has constant one.
\end{lemma}

\begin{proof}
Put $\Delta_0^i=P_0^i$ and $\Delta_n^i=P_n^i-P_{n-1}^i$ for $n\ge1$.
Within each family these are pairwise orthogonal projections.  A prefix of
the index set $\{0,\ldots,T\}$ can be partitioned into dyadic index blocks,
with at most one block of each length.  For a block $A$ put
$\Delta_A^i=\sum_{n\in A}\Delta_n^i$.  At any fixed block length, the
operators $\Delta_A^i$ are pairwise orthogonal projections.

Expand each factor of $Z_{n(\nu)}$ using this binary decomposition.  Fix a
tuple of block lengths, one for each of the $m$ factors.  A term is zero if
one of those lengths does not occur in the relevant prefix.  Otherwise it
has the form
\[
 B_\nu\Delta_{A_m(\nu)}^m\cdots\Delta_{A_1(\nu)}^1g.
\]
Group rows with the same tuple of blocks.  The row contraction bounds their
sum of squared norms by
$\|\Delta_{A_m}^m\cdots\Delta_{A_1}^1g\|^2$.  Summing over the block
tuples that occur is bounded by the sum over all tuples of those fixed
lengths.  Successive Bessel inequalities give
\[
 \sum_{A_1,\ldots,A_m}
       \|\Delta_{A_m}^m\cdots\Delta_{A_1}^1g\|^2\le\|g\|^2.
\]
This argument uses orthogonality within each layer, not commutation between
layers.  There are at most $L_T^m$ tuples of block lengths.  The triangle
inequality in the Hilbert direct sum of the row spaces proves the first
assertion.

For the second, telescope in the order
\[
 Z_n-Z_{n-1}
   =\sum_{i=1}^m
       P_n^m\cdots P_n^{i+1}\Delta_n^i
                       P_{n-1}^{i-1}\cdots P_{n-1}^1.
\]
Discard the left contractions and use Cauchy--Schwarz in the sum over $i$.
For a fixed $i$, the rows $B_n=\Delta_n^i$, $1\le n\le T$, satisfy the row
contraction hypothesis.  Apply the first assertion to the $i-1$ preceding
families and the prefix index $n-1$.  The resulting bound is
\[
 \sum_{n=1}^T\|(Z_n-Z_{n-1})g\|^2
    \le m\sum_{i=1}^m L_T^{2(i-1)}\|g\|^2
    \le m^2L_T^{2(m-1)}\|g\|^2,
\]
as claimed.
\end{proof}

\subsection{Compression of local projection products}

\begin{lemma}[Compression]
\label{lem:compression}
Use the column system and budget $D\ge2$ above on a finite dyadic tree with
root $J$.  Let $P^1,\ldots,P^m$ be inherited mask families, with their local
orthogonal projections.  Suppose finitely many bounded linear rows
\[
 B_\nu:L^2(J;\mathcal H)\longrightarrow\mathcal Y_\nu
\]
satisfy
\[
 \sum_\nu\|B_\nu z\|_{\mathcal Y_\nu}^2\le\|z\|_{L^2(J;\mathcal H)}^2,
 \qquad B_\nu=B_\nu\1_{u(\nu)},
\]
where each $u(\nu)$ is a cell of the tree.  Then an absolute constant $C$
satisfies
\[
 \left(\sum_\nu
    \|B_\nu P_{u(\nu)}^m\cdots P_{u(\nu)}^1z\|^2\right)^{1/2}
 \le \bigl(C(m+1)^C\log(2D)\bigr)^{C(m+1)}\|z\|_{L^2(J;\mathcal H)}.
\]
The estimate is uniform in the row cells, the total number of labels, and the
number of spatial depths.  Identity factors may be omitted.  Products
containing complements $1-P^i$ satisfy the same form of bound after changing
$C$.  Integrating such compositions against a fixed finite measure costs, in
addition, at most the total variation of that measure.
\end{lemma}

\begin{proof}
For a nonempty product, we first divide the row cells into size segments
and, on each segment, replace the input by a bounded function with the
same first-layer moments.  We then stop when the last-layer coefficient
vector moves: nested-prefix estimates control the anchor products, and
the pathwise birth-label budget controls the coefficient errors.
Attempts are summed using both length and energy contraction.  Finally,
the size segments are summed using the bounds for their weighted
averages and representative energies.

The assertion for an empty product is the row contraction hypothesis.  Assume
$m\ge1$, and abbreviate the ordinary $L^2(J;\mathcal H)$ norm to $\|\cdot\|_2$.
Add one final dyadic generation below all row cells, retaining columns and
refitting there if necessary.  This changes neither the hypothesis nor the
quantity to be estimated.

\paragraph{Size segments and a bounded moment representative.}
For every cell $K$ put
\[
 a_K=\fint_K\|z(y)\|_{\mathcal H}\,dy.
\]
Starting from $J$, stop at the first proper descendant cells for which the
average is greater than $2a_J$, and repeat inside each stopping cell.  It is
enough to run this stopping rule through the last possible row depth.  A root
with zero average contributes no rows and can be omitted.  Denote the
resulting stopping roots by $\mathcal S$.

Assign a row cell $u$ to the deepest root $K\in\mathcal S$ containing it.
More generally, let $\mathcal T_K$ be the full subtree of cells in $K$ that
are not contained in a proper stopping child of $K$, through the last row
depth.  Thus every assigned row cell lies in $\mathcal T_K$.  Cut $K$ at its
next stopping roots and at final leaves on all other paths; denote this
partition by $\pi_K$.  Every assigned row cell is partitioned by cells of
$\pi_K$.  On a cut cell $Q$ the average of $\|z\|$ is at most $4a_K$: its
dyadic parent has average at most $2a_K$, and a child average is at most twice
its parent's average.

For $Q\in\pi_K$, take all first-layer labels used by cells of
$\mathcal T_K$ containing $Q$, and project $z|_Q$ onto their restrictions to
$Q$.  Let $g_K$ be the direct sum of these projections, extended by zero off
$K$.  At most $D$ labels occur on each such path.  By
Lemma~\ref{lem:private-columns}, with
\[
 G=4D^5,
\]
we have
\[
 \|g_K\|_{L^\infty(K;\mathcal H)}\le Ga_K,
 \qquad \|g_K\|_2^2=\langle z,g_K\rangle.
\]
The second identity follows from orthogonality on each cut cell.  If
$u\in\mathcal T_K$ and $v\in\mathcal V^1(u)$, then $v$ was included on
every cut cell inside $u$.  Adding its moment identities over that partition
of $u$ gives
\[
 P_u^1g_K=P_u^1z,
 \qquad P_u^m\cdots P_u^1g_K=P_u^m\cdots P_u^1z.
\]

Let $M_{\mathrm d}\|z\|$ be the dyadic maximal function within $J$.
Stopping averages more than double on every path.  Consequently
\[
 \sum_{K\in\mathcal S}a_K^2|K|
    \le\frac43\|M_{\mathrm d}\|z\|\|_2^2
    \le\frac{16}{3}\|z\|_2^2.
\]
We also need the stronger bound for the moment representatives
\[
 S:=\sum_{K\in\mathcal S}\|g_K\|_2^2
       \le C\log(2D)\|z\|_2^2.
\]
Here is a proof that avoids any factor which is a power of $D$.  At a fixed
point $y$, split the roots containing $y$ according as
\[
 a_K<\frac{M_{\mathrm d}\|z\|(y)}{GD^2}
 \quad\text{or}\quad
 a_K\ge\frac{M_{\mathrm d}\|z\|(y)}{GD^2}.
\]
The sum of $a_K$ in the first class is at most twice its displayed upper
threshold.  Thus their total pointwise contribution to
$\sum_K\|g_K(y)\|$ is at most $2D^{-2}M_{\mathrm d}\|z\|(y)$.  The second
class has at most
\[
 L_D=2+\left\lceil\log_2(GD^2)\right\rceil\le C\log(2D)
\]
members, since $a_K\le M_{\mathrm d}\|z\|(y)$ and the averages double.
Using the moment identity, pointwise Cauchy--Schwarz for the second class,
and then Cauchy--Schwarz in $y$, we obtain
\[
 S\le 2D^{-2}\|z\|_2\|M_{\mathrm d}\|z\|\|_2
          +\sqrt{L_D}\,\|z\|_2\sqrt S
   \le4D^{-2}\|z\|_2^2+\sqrt{L_D}\,\|z\|_2\sqrt S.
\]
Solving this quadratic inequality in $\sqrt S$ proves the asserted bound.

\paragraph{Anchors within one size segment.}
Fix $K\in\mathcal S$, write $g=g_K$ and $a=a_K$, and put
\[
 \tau=D^{-10},\qquad
 b=2D(m+1),\qquad
 T=\left\lceil b^{128(m+1)^2}\right\rceil.
\]
For each cell $u\in\mathcal T_K$, let $c(u)$ be the unique last-layer
coefficient vector of $P_u^m\cdots P_u^1g$, extended by zero to all labels.

We describe an attempt from a cell $R\in\mathcal T_K$, initially $R=K$.
Declare $R$ an anchor.  From an anchor $A$, choose as its next anchors the
first proper descendant cells $u\in\mathcal T_K$ for which
\[
 \|c(u)-c(A)\|_{\ell^2(\mathcal V)}>\tau a.
\]
Stop a path at its $T$th transition and send that cell and its descendants
to a new attempt, using the same $g$ and $a$.  Denote the disjoint new roots
from $R$ by $\mathcal R(R)$.

To apply Lemma~\ref{lem:projection-prefixes}, complete each anchor generation
to a stopping partition $\pi_n$ of $R$, $0\le n\le T$.  The active atoms are
the anchors of generation $n$.  The other atoms are terminal cut cells on
paths that have ended; leave those atoms unchanged in later partitions.
When refining an active anchor, its next anchors together with the terminal
cut cells not lying inside a next anchor partition it.  Use the original
masks on active anchors.  On a terminal cut one can use an inherited local
refit, for example with its original mask if that mask is already defined,
and keep this local projection unchanged afterward.  Hence these are
refining partitions, and the resulting direct-sum auxiliary projection
families $P_n^i=P_{\pi_n}^i$ are nested.  Put
$Z_n=P_n^m\cdots P_n^1$ on $L^2(R;\mathcal H)$.

For any $R'\in\mathcal R(R)$, its ancestors in the anchor tree have
generations $0,\ldots,T$.  On $R'$ the restriction of
$(Z_n-Z_{n-1})(\1_Rg)$ is the difference of the constant-coefficient
combinations at the corresponding two consecutive anchors.  Inheritance
makes all the older labels available at the newer anchor.  The private
coordinate bound in Lemma~\ref{lem:private-columns} therefore gives
\[
 \|(Z_n-Z_{n-1})(\1_Rg)(y)\|_{\mathcal H}
         >D^{-1}\tau a
 \quad\text{for almost every }y\in R',\quad 1\le n\le T.
\]
Integrating over the disjoint new roots, summing in $n$, and applying
Lemma~\ref{lem:projection-prefixes} yields
\[
 T D^{-2}\tau^2a^2\sum_{R'\in\mathcal R(R)}|R'|
        \le m^2L_T^{2m}\|\1_Rg\|_2^2.
\]
Since $\|g\|_\infty\le Ga$, this implies both
\[
 \sum_{R'\in\mathcal R(R)}|R'|\le E|R|,
 \qquad
 \sum_{R'\in\mathcal R(R)}\|\1_{R'}g\|_2^2
       \le E\|\1_Rg\|_2^2,
 \qquad
 E=\frac{16D^{32}m^2L_T^{2m}}{T}.
\]
Our fixed choice of $T$ makes $E<1/2$.  For completeness, $b\ge8$ and
\[
 L_T\le129(m+1)^2\log_2 b\le b^4.
\]
The second inequality follows from $m+1\le b/4$ and
$\log_2b\le b/2$.  Thus
\[
 E\le16b^{34+8m-128(m+1)^2}
       \le16\,8^{-470}<\tfrac12.
\]
It follows by generation-by-generation summation that all attempt roots
inside this size segment, denoted by $\mathcal R_K$, satisfy
\[
 \sum_{R\in\mathcal R_K}|R|\le2|K|,
 \qquad
 \sum_{R\in\mathcal R_K}\|\1_Rg_K\|_2^2\le2\|g_K\|_2^2.
\]

\paragraph{Rows assigned to an attempt.}
Assign a row in the size segment to the deepest attempt root containing its
cell.  For a row cell $u$ assigned to the attempt from $R$, let $A$ be its
last ancestor anchor in that attempt, and let $j$ be the generation of $A$.
Here $0\le j<T$.  The cell $A$ is an atom of $\pi_j$ and contains $u$.
Consequently the restriction of $Z_j(\1_Rg_K)$ to $u$ is exactly the product
at the anchor, restricted to $u$.  The stopping rule gives
\[
 \|c(u)-c(A)\|_{\ell^2(\mathcal V)}\le\tau a_K.
\]
The support condition on $B_\nu$ now permits the decomposition of its input
product into this anchor product and its coefficient error.  The rows
assigned to $R$, restricted to $L^2(R;\mathcal H)$, still satisfy the row
contraction hypothesis.  Lemma~\ref{lem:projection-prefixes} shows that their
anchor contributions have squared sum at most
\[
 L_T^{2m}\|\1_Rg_K\|_2^2.
\]
This is the step in which arbitrary, unequal physical depths of the row
cells cause no additional cost: the prefix index is the row's last anchor
generation, and its cell lies wholly in that anchor atom.

For the error write $e_{\nu,v}=c(u(\nu))_v-c(A(\nu))_v$.  It is supported
on labels available on $u(\nu)$ and has squared coefficient sum at most
$\tau^2a_K^2$.  Since $u(\nu)\subset R\cap J_v$ whenever this coefficient
can be nonzero, locality of $B_\nu$ gives
\[
 B_\nu\left(\1_{u(\nu)}\sum_ve_{\nu,v}\Phi_v\right)
       =\sum_ve_{\nu,v}B_\nu(\1_R\widetilde\Phi_v).
\]
Cauchy--Schwarz in $v$, followed by the row contraction for each fixed
label, bounds the squared sum of errors in this attempt by
\begin{align*}
 &\tau^2a_K^2
     \sum_v\sum_{\substack{\nu\text{ assigned to }R\\
                            v\text{ available on }u(\nu)}}
           \|B_\nu(\1_R\widetilde\Phi_v)\|^2 \\
 &\hspace{20mm}\le
      \tau^2a_K^2\sum_v\|\1_R\widetilde\Phi_v\|_2^2
      \le\tau^2a_K^2D^3|R|.
\end{align*}
The last inequality uses the pathwise budget in its integral form:
\[
 \sum_v\|\1_R\widetilde\Phi_v\|_2^2
   =\int_R\sum_{v:y\in J_v}\|\Phi_v(y)\|_{\mathcal H}^2\,dy
   \le D^3|R|.
\]
In particular, no estimate of the total number of labels on different
branches has been used.

\paragraph{Summation and the final constant.}
Each row belongs to exactly one size segment and one attempt within it.
The triangle inequality in the full row Hilbert space, the two packing
bounds, and the estimates for $S$ and the stopping averages give
\begin{align*}
 \left(\sum_\nu
     \|B_\nu P_{u(\nu)}^m\cdots P_{u(\nu)}^1z\|^2\right)^{1/2}
 &\le
    L_T^m\left(2\sum_K\|g_K\|_2^2\right)^{1/2}
    +\tau D^{3/2}
          \left(2\sum_Ka_K^2|K|\right)^{1/2} \\
 &\le C\left(L_T^m\sqrt{\log(2D)}+D^{-17/2}\right)\|z\|_2.
\end{align*}
Finally,
\[
 L_T\le C(m+1)^2\log(2D(m+1))
       \le C(m+1)^3\log(2D),
\]
which is bounded by the stated expression after increasing its absolute
constant.  A complement is expanded as $1-P^i$; the at most $2^m$ resulting
ordered products obey the same form of bound, and their number is absorbed
in that constant.  For a fixed average of compositions, apply the estimate
to each member and use the integral triangle inequality in the row Hilbert
space.  This proves the remaining assertions.
\end{proof}

\subsection{Rational heights and occupied scale ranges}

The next observation is useful when real frequencies are unrestricted but
all coefficients used to combine them have bounded rational height.  Small
divisors can place the relevant scale ranges far apart; they do not create
arbitrarily many such ranges.

\begin{lemma}[Height bins]
\label{lem:height-bins}
Let $d\ge1$, $B\ge2$, and $\xi\in\R^d$.  Let
\[
 \mathcal E=\left\{|a\cdot\xi|:
       a\in\Q^d\text{ has height at most }B,\ a\cdot\xi\ne0\right\}.
\]
There is an absolute constant $C$ such that the following assertions hold.
\begin{enumerate}
\item The set $\{\log_2v:v\in\mathcal E\}$ is contained in a union of at
most $d$ intervals, each of length at most
$C(d+1)^3\log(2B)$.  In particular, $\mathcal E$ meets at most
$C(d+1)^4\log(2B)$ dyadic magnitude bins.
\item If $r_n=r_0 2^{-n}$ with $r_0>0$, $\ell\in\{1,2\}$, and
$0<\alpha\le\beta$, then
\[
 \#\{n\in\Z:\text{ some }v\in\mathcal E
                  \text{ satisfies }\alpha\le r_n^\ell v\le\beta\}
 \le C(d+1)^4\bigl(\log(2B)+\log(2\beta/\alpha)\bigr).
\]
\item For every $\varepsilon>0$, with the maximum of an empty set defined
to be zero,
\[
 \sum_{n\in\Z}
     \max\bigl(\{r_n^\ell v:v\in\mathcal E,
                       r_n^\ell v\le\varepsilon\}\cup\{0\}\bigr)
     \le C(d+1)^4\log(2B)\,\varepsilon.
\]
\end{enumerate}
The same conclusions apply to suprema of finitely many such coordinate
evaluations.  For Euclidean norms of evaluation vectors with at most $h$
coordinates, one can add $C d\log(2h)$ to the bounds for the number of
occupied bins and scales, and to the coefficient of $\varepsilon$ in the
last estimate.  No bound on the size of $\xi$ is required.
\end{lemma}

\begin{proof}
If $\mathcal E$ is empty there is nothing to prove.  It is otherwise a
finite set.  We first bound coefficients in a minimal rational dependence.
Suppose $a_1,\ldots,a_k$ are minimally dependent vectors of the indicated
height.  Then $k\le d+1$ and their rank is $k-1$.  For each $i$, clear the
denominators in its $d$ entries by a common denominator $q_i\le B^d$.
The integer vector $v_i=q_i a_i$ has each entry bounded by $B^d$.  Select
$k-1$ coordinate rows on which the vectors have rank $k-1$.  Their signed
maximal minors give a null vector $(b_1,\ldots,b_k)$ with
\[
 0<|b_i|\le d! B^{d^2}.
\]
None of these minors is zero: a zero component of a null vector would give
a dependence on a proper subset, contrary to minimality.  Returning to
the original vectors gives
\[
 \sum_{i=1}^k c_i a_i=0,
 \qquad c_i=b_iq_i\in\Z\setminus\{0\},
 \qquad |c_i|\le H:=d! B^{d^2+d}.
\]

Set $R=4dH$.  Suppose coefficient vectors have nonzero evaluation magnitudes
\[
 0<s_1<s_2<\cdots<s_k,
 \qquad s_{i+1}>Rs_i.
\]
If these vectors were dependent, apply the preceding construction to a
minimal dependent subset, ordered by evaluation magnitude, and let $s_*$
be its largest magnitude.  The coefficient of that term has absolute value
at least one, whereas all the other terms together have absolute value at
most
\[
 dH\frac{s_*}{R}<s_*.
\]
They cannot sum to zero.  Thus vectors whose evaluation magnitudes are
successively separated by a factor greater than $R$ are linearly
independent, and there can be at most $d$ of them.

Choose the smallest element $s$ of $\mathcal E$, cover $[s,Rs]$, and repeat
with the smallest still uncovered element.  The selected elements are
successively separated by more than $R$, so at most $d$ intervals are
needed.  Their logarithmic lengths are $\log_2R$, and
\[
 \log_2R
    \le C(d+1)^2\log(2dB)
    \le C(d+1)^3\log(2B).
\]
Each logarithmic interval meets at most $2+\log_2R$ unit bins.  This proves
the first assertion.

For the second, consider one logarithmic value interval $[s,s+\log_2R]$.
The condition $\alpha\le r_n^\ell v\le\beta$ confines $n$ to a real
interval of length at most
\[
 \frac{\log_2R+\log_2(\beta/\alpha)}{\ell}.
\]
That interval has at most two plus its length integer points.  Summing over
the at most $d$ value intervals and using the preceding bounds proves the
claimed scale count.  The factor $r_0$ translates the interval and has no
effect on its length.

For the third, partition $\mathcal E$ by its occupied bins
$[2^j,2^{j+1})$.  For one such bin, a contribution at scale $n$ is possible
only if $r_n^\ell2^j\le\varepsilon$, and its maximum is at most
$r_n^\ell2^{j+1}$.  Let $n_0$ be the first integer for which the former
inequality holds.  The sum of these upper bounds for $n\ge n_0$ is at most
\[
 2\varepsilon\sum_{a=0}^{\infty}2^{-\ell a}
       \le4\varepsilon.
\]
The maximum over all occupied bins is bounded by the sum of their separate
maxima.  Multiplying by the occupied-bin count proves the assertion.

Finally, the supremum norm of a nonzero evaluation vector is the magnitude
of one of its nonzero coordinates, so the same set $\mathcal E$ contains
all such suprema.  If there are at most $h$ coordinates, its Euclidean norm
lies between that supremum and $\sqrt h$ times the supremum.  Enlarging each
of the $d$ logarithmic intervals to the right by
$\tfrac12\log_2h$ covers all these Euclidean norms.  The preceding bin,
scale, and geometric-tail arguments then give exactly the additional costs
stated in Lemma~\ref{lem:height-bins}.
\end{proof}

In particular, these estimates permit the rational coefficient vector to
be chosen anew at each depth: every such choice still lies in the same
bounded-height family.  They also apply after adjoining finitely many fixed
real parameters to $\xi$, for example the coordinates of a horizontal
frequency vector and their pairwise products.
